\documentclass[../main.tex]{subfiles}
\begin{document}

\chapter{マルチプロセッシング}
\lessoninfo{
  video={Lesson 18，動画 1--22},
  textbook={H\&P \cite{hennessy2017} §1.1--1.2，§1.5，§5.1，第 3 章，付録 B},
  keywords={Flynn の分類，MIMD，マルチコアプロセッサ，集中共有メモリ，UMA，SMP，分散共有メモリ，NUMA，メッセージパッシング，共有メモリ，チップマルチプロセッサ，細粒度マルチスレッディング，同時マルチスレッディング，ワーキングセット},
}

ここまでの章では，パイプライン化，分岐予測，アウトオブオーダ実行，
メモリ階層によって，1つのプロセッサが1つのプログラムをできる限り速く実行
できるようにしてきた．そして第17章では，そのようなシステムを信頼できる
ものにした．この章では，複数のプログラム，あるいは1つのプログラムの
複数の部分を同時に実行するマシンを扱う．並列マシンを分類し，各技術
世代がもたらす追加のトランジスタが，今日ではなぜより速い単一コアでは
なく追加のコアになるのかを説明し，プロセッサをメモリに接続する2つの方式と
それぞれに対応する2つのプログラミングモデルを比較し，最後にマルチ
スレッディングのハードウェア---とりわけ，1つのコアに複数のスレッドを
同時に実行させる同時マルチスレッディング---で締めくくる．

\section{Flynn の分類}
\label{sec:l18-flynn}

計算機が一度に複数のことを行う方法には，いくつかの種類がある．
\source{Lesson 18，動画 1--2；H\&P §1.2；\cite{flynn1966}．}並列マシンの
古典的な分類である\term[ふりんのぶんるい]{Flynn の分類}[Flynn's taxonomy]%
は，2つの数に着目する．マシンが同時に実行する\emph{命令ストリーム}%
---それぞれが独自のプログラムカウンタを持つ命令の列---の数と，それらの
命令が操作する\emph{データストリーム}の数である．どちらの数についても
「1つ」と「2つ以上」を区別すると，4つのクラスが得られる
（\cref{tab:l18-flynn}）．

\begin{table}[htbp]
  \centering\small
  \caption{並列マシンの Flynn の分類．}
  \label{tab:l18-flynn}
  \begin{tabular}{@{}lccl@{}}
    \toprule
    クラス & 命令ストリーム & データストリーム & 例 \\
    \midrule
    SISD & 1        & 1        & 単一プロセッサ \\
    SIMD & 1        & $>1$     & ベクトルプロセッサ，SSE/AVX \\
    MISD & $>1$     & 1        & ストリームプロセッサ（稀） \\
    MIMD & $>1$     & $>1$     & マルチプロセッサ，マルチコア \\
    \bottomrule
  \end{tabular}
\end{table}

\begin{description}
  \item[SISD] 1つの命令ストリームを1つのデータストリームに対して実行する
    マシンが \term{SISD}[single instruction, single data] である．これまで
    の章の単一プロセッサがこれである．スーパースカラのアウトオブオーダコアも
    依然として SISD である．同じサイクルに実行される命令は，すべて1つの
    プログラムに属するからである．
  \item[SIMD] 1つの命令ストリームを実行するが，各命令が複数のデータ要素
    に対して演算するマシンが
    \term{SIMD}[single instruction, multiple data] である．ベクトル
    プロセッサがこの種であり，現在の命令セットの
    マルチメディア拡張（x86 の MMX，SSE，AVX）も同様で，その命令は短い
    ベクトルのすべての要素に1つの演算を適用する．
  \item[MISD] 複数の命令ストリームを1つのデータストリームに対して実行
    するマシンが \term{MISD}[multiple instruction, single data] である．
    複数のプログラムがそれぞれ同じ入力データを処理するストリームプロセッサが
    その例になる．この種のマシンは稀である．
  \item[MIMD] 複数の命令ストリームを，それぞれ独自のデータに対して実行
    するマシンが \term{MIMD}[multiple instruction, multiple data] である．
\end{description}

複数のプロセッサからなり，各プロセッサが自身のデータに対して自身の
プログラム（あるいは1つのプログラムの自身の担当部分）を実行するマシンを
\term[まるちぷろせっさ]{マルチプロセッサ}[multiprocessor]という．
マルチプロセッサは MIMD マシンであり，今日使われている並列マシンの
ほとんどすべてがこのクラスに属する．

\section{なぜマルチプロセッサか}
\label{sec:l18-why}

単一プロセッサを速くする方法は2つある．サイクル当たりにより多くの命令を
完了するか，サイクルをより短くするかである．\source{Lesson 18，動画
3--4；H\&P §1.1，§1.5．}どちらも余地を使い果たしている．
\begin{itemize}
  \item \textbf{幅を広げる．} 現在のコアはすでにサイクル当たり約4命令を
    扱う．プログラムがそれ以上のスロットを埋めるだけの独立な命令を含む
    ことは稀であり，幅をさらに増やしても収穫逓減となる（第6章）．
  \item \textbf{速くする．} クロック周波数を上げるには電源電圧を上げる
    必要がある．\margin{代わりにクロック周波数が上がらなくなった．
    Pentium~4 は 2004 年に \SI{3.8}{\giga\hertz} に達し，今日の最も速い
    デスクトップ向け製品でもおよそ \SI{6}{\giga\hertz} である．H\&P
    §1.5．}%
    動的電力は $V^{2}f$ に比例する
    （\cref{eq:active-power}）ので，$f$ と $V$ をともに上げると電力は
    おおよそ $f^{3}$ で増える．しかもプロセッサはすでに，妥当なコストで
    除去できるだけの電力を消費している．
\end{itemize}

一方でムーアの法則は続いている．1つのチップに収まるトランジスタ数は，
今なおおよそ18か月で倍増する．より複雑な単一コアを作る代わりに，その
トランジスタでコアを複製することができる．
\term[まるちこあぷろせっさ]{マルチコアプロセッサ}[multicore processor]%
とは，2個以上の完全なプロセッサ---\emph{コア}と呼ばれる---を1つのチップ
に収めたものである．世代ごとにコア数を倍にすれば，世代ごとに性能を倍に
できる---ソフトウェアがすべてのコアを使う限りにおいて．\caution{ムーア
自身の数字は 1965 年には1年\cite{moore1965}，1975 年の改訂では2年で
ある．「18か月」は後年の通俗版である．}

\begin{example}
  \label{ex:l18-power}
  あるコアが \SI{3}{\giga\hertz} で動作し，\SI{40}{\watt} の動的電力を
  消費する．これを $1.2$ 倍速い \SI{3.6}{\giga\hertz} で動作させるには，
  電圧を比例して上げる必要があり，電力はおよそ $1.2^{3}\approx 1.73$ 倍
  の \SI{69}{\watt} になる．元のコアを \SI{3}{\giga\hertz} で2個動かすと
  \SI{80}{\watt} を消費し，1個のコアの2倍の仕事を---プログラムが両方を
  忙しく保てば---完了する．\SI{2.4}{\giga\hertz}（$0.8$ 倍）で2個動かせば
  $2\times\SI{40}{\watt}\times 0.8^{3}\approx\SI{41}{\watt}$ となり，元の
  単一コア1個とほぼ同じ消費でありながら，合わせてその $1.6$ 倍の仕事を
  完了する．
\end{example}

\section{マルチプロセッサには並列プログラムが必要である}
\label{sec:l18-parallel-programs}

単一プロセッサを速くすれば，遠い昔に書かれたものも含め，あらゆる
プログラムが速くなる．マルチプロセッサが速くするのは，仕事を複数のコアに
分割するプログラムだけである．\source{Lesson 18，動画 5；H\&P §5.1．}%
逐次の（シングルスレッドの）コードと比べて，そうした並列プログラムは
あらゆる面でより難しい．
\begin{itemize}
  \item \textbf{開発．} 問題を解くことに加えて，プログラマは仕事を
    スレッドに分割し，スレッド間を協調させなければならない．
  \item \textbf{デバッグ．} 並列プログラムの振る舞いはスレッド間の相対的
    なタイミングに依存し，それは実行のたびに変わるので，ある実行で現れた
    誤りが次の実行では現れないことがある．
  \item \textbf{性能のスケーリング．} 理想的には，性能はコア数に比例して
    伸びる．実際にはある程度のコア数までは伸び，その後は頭打ちになる．
    プログラムを改良すればその水準は上がるが，それでも性能はどこかで
    伸びなくなる（\cref{fig:l18-scaling}）．
\end{itemize}

\begin{marginfigure}
  \resizebox{\linewidth}{!}{% Speedup versus number of cores: ideal linear scaling, and the levelling
% off predicted by Amdahl's law for parallel fractions 0.97 and 0.9.
\begin{tikzpicture}[x=0.118cm,y=0.118cm, font=\scriptsize\sffamily]
  \draw[ln-rule] (0,0) grid[xstep=8,ystep=8] (32,32);
  \draw[->] (0,0) -- (34,0);
  \draw[->] (0,0) -- (0,34);
  \foreach \v in {0,8,16,24,32} {
    \node[below, inner sep=1.5pt] at (\v,0) {\v};
    \node[left, inner sep=1.5pt] at (0,\v) {\v};
  }
  \node[below] at (16,-3.2) {\iflangja{コア数 $N$}{cores $N$}};
  \node[rotate=90, above] at (-4.6,16) {\iflangja{高速化率}{speedup}};
  \draw[thick, ln-muted, dashed] (1,1) -- (32,32)
    node[pos=0.78, above left, inner sep=1pt] {\iflangja{理想}{ideal}};
  \draw[thick, ln-accent] plot[domain=1:32, samples=48, smooth]
    (\x, {1/(0.03+0.97/\x)});
  \node[ln-accent, above left, inner sep=1pt] at (32,16.6) {$f_{\mathrm{par}}=0.97$};
  \draw[thick, ln-caution] plot[domain=1:32, samples=48, smooth]
    (\x, {1/(0.1+0.9/\x)});
  \node[ln-caution, above left, inner sep=1pt] at (32,7.8) {$f_{\mathrm{par}}=0.9$};
\end{tikzpicture}
}
  \caption{$N$ コアでの高速化率: 理想的なスケーリングと
    \cref{eq:l18-amdahl-cores}．}
  \label{fig:l18-scaling}
\end{marginfigure}
頭打ちが避けられない理由は Amdahl の法則（\cref{eq:amdahl}）が示す．
プログラムの実行時間のうち割合 $f_{\mathrm{par}}$ が $N$ 個のコアに均等に
分割され，残りが逐次のままであるとすると，高速化率は高々
\begin{equation}
  S(N) \;=\; \frac{1}{(1-f_{\mathrm{par}}) + \dfrac{f_{\mathrm{par}}}{N}}
  \;<\; \frac{1}{1-f_{\mathrm{par}}}
  \label{eq:l18-amdahl-cores}
\end{equation}
である．スレッド間の協調にはこれに加えて時間がかかるので，実際の
プログラムはこれらの曲線より下に留まる．\tip{理想の半分の高速化率
$S(N)=N/2$ には $f_{\mathrm{par}}=(N-2)/(N-1)$ が必要である．8 コアで
\SI{86}{\percent}，64 コアで \SI{98}{\percent} である．}

\begin{example}
  あるプログラムでは，時間の \SI{95}{\percent} が，反復が互いに独立な
  ループに費やされており，$f_{\mathrm{par}}=0.95$ である．8 コアでの
  高速化率は高々 $1/(0.05+0.95/8)\approx 5.9$，32 コアでは
  $1/(0.05+0.95/32)\approx 12.5$ であり，コアをいくら増やしても
  $1/0.05 = 20$ を超えることはない．
\end{example}

\needspace{5\baselineskip}
\begin{keypoint}
  コアを複製することは，より大きな，あるいはより速い単一コアを作るより
  はるかに効率よくトランジスタをスループットに変える．ただしそれは，
  多数のコアを使うように書かれたプログラムに対してだけであり，しかもその
  逐次部分が許す範囲までである．
\end{keypoint}

\section{集中共有メモリ}
\label{sec:l18-centralized}

最も単純なマルチプロセッサは，すべてのプロセッサを1つの主記憶に接続する．
\source{Lesson 18，動画 6--8；H\&P §5.1．}

\begin{definition}[集中共有メモリ]
  \term[しゅうちゅうきょうゆうめもり]{集中共有メモリ}[centralized shared memory]%
  のマルチプロセッサでは，それぞれ自身のキャッシュを持つすべてのコアが，
  1本のバスによって1つの主記憶と I/O 装置に接続される
  （\cref{fig:l18-centralized}）．
\end{definition}

\begin{figure}[htbp]
  \centering
  % Centralized shared memory (UMA, SMP): cores with private caches share one
% main memory and the I/O devices through a common bus.
\begin{tikzpicture}[x=1cm,y=1cm, font=\footnotesize\sffamily,
  core/.style={draw, fill=ln-accent!15, minimum width=1.7cm, minimum height=0.6cm, inner sep=2pt},
  cache/.style={draw, fill=ln-box, minimum width=1.7cm, minimum height=0.5cm, inner sep=2pt},
  big/.style={draw, fill=ln-box, minimum height=0.7cm, inner sep=2pt}]
  \foreach \i in {0,...,3} {
    \node[core] (c\i) at (\i*2.3,0) {\iflangja{コア}{core}~\i};
    \node[cache] (k\i) at (\i*2.3,-0.85) {\iflangja{キャッシュ}{cache}};
    \draw (c\i.south) -- (k\i.north);
    \draw (k\i.south) -- (\i*2.3,-1.7);
  }
  \draw[line width=1.4pt] (-1.0,-1.7) -- (7.9,-1.7)
    node[right, font=\footnotesize] {\iflangja{バス}{bus}};
  \node[big, minimum width=4.4cm] (mem) at (1.15,-2.75) {\iflangja{主記憶}{main memory}};
  \node[big, minimum width=2.8cm] (io) at (5.75,-2.75) {\iflangja{I/O 装置}{I/O devices}};
  \draw (mem.north) -- (1.15,-1.7);
  \draw (io.north) -- (5.75,-1.7);
\end{tikzpicture}

  \caption{集中共有メモリ: すべてのコアが同じバスを通じて主記憶と I/O
    装置に到達する．}
  \label{fig:l18-centralized}
\end{figure}

どのコアも任意のメモリ位置を読み書きでき，任意の I/O 装置を使えるので，
コア間のデータ交換は，一方がメモリに書き，他方がそれを読むだけで済む．
メモリはすべてのコアから等しく離れているため，どのコアもある位置への
アクセスに同じ時間を要する．すなわちこのマシンは
\term{UMA}[uniform memory access] である．また，
どのコアも特別な役割を持たないため，このマシンは
\term[たいしょうがたまるちぷろせっさ]{対称型マルチプロセッサ}[symmetric multiprocessor]%
\index{SMP|see{対称型マルチプロセッサ}}（SMP）とも呼ばれる．チップに
接続された主記憶をコアが共有するマルチコアチップが，最も一般的な例で
ある．

集中共有メモリは，2つの理由から多数のコアにはスケールしない．
\begin{itemize}
  \item \textbf{メモリ容量．} 1つのメモリがすべてのコアのデータを保持
    するので，コア数とともに大きくしなければならず，メモリは大きいほど
    遅い（第12章と第15章）．
  \item \textbf{メモリ帯域．} すべてのコアのキャッシュミスと，すべての
    I/O 転送が同じバスを通る．コアが増えるほど互いのメモリアクセスを待つ
    ようになり，ついにはメモリアクセスが事実上直列化され，コアを追加して
    も利益が得られなくなる．
\end{itemize}
したがって集中共有メモリがうまく働くのは，コア数がおよそ 2 から 16 の
小規模なマシンに限られる．

\section{分散共有メモリ}
\label{sec:l18-distributed}

より大規模なマシンは，メモリをプロセッサに分散させる．
\source{Lesson 18，動画 9--10；H\&P §5.1．}

\begin{definition}[分散共有メモリ]
  \term[ぶんさんきょうゆうめもり]{分散共有メモリ}[distributed shared memory]%
  のマルチプロセッサでは，各コアが自身のローカルメモリとネットワーク
  インタフェースを持ち，コアどうしはネットワークで接続される
  （\cref{fig:l18-distributed}）．コアが直接アクセスするのは自身の
  ローカルメモリだけであり，他のコアが保持するデータは，そのコアに
  ネットワーク経由で要求を送って得る．
\end{definition}

\begin{figure}[htbp]
  \centering
  % Distributed shared memory (NUMA): every node has its own memory and a
% network interface; memory of other nodes is reached through the network.
\begin{tikzpicture}[x=1cm,y=1cm, font=\footnotesize\sffamily, >={Stealth[length=1.8mm]},
  core/.style={draw, fill=ln-accent!15, minimum width=2.1cm, minimum height=0.55cm, inner sep=1pt},
  part/.style={draw, fill=ln-box, minimum width=2.1cm, minimum height=0.5cm, inner sep=1pt},
  small/.style={draw, fill=ln-box, minimum width=0.95cm, minimum height=0.5cm, inner sep=1pt, font=\scriptsize\sffamily}]
  \foreach \i in {0,...,3} {
    \begin{scope}[xshift=\i*2.75cm]
      \draw[ln-rule, rounded corners=2pt] (-1.25,0.5) rectangle (1.25,-2.75);
      \node[core] (c\i) at (0,0) {\iflangja{コア}{core}~\i};
      \node[part] (k\i) at (0,-0.8) {\iflangja{キャッシュ}{cache}};
      \node[small] (m\i) at (-0.575,-2.2) {\iflangja{メモリ}{memory}};
      \node[small] (n\i) at (0.575,-2.2) {NIC};
      \draw (c\i.south) -- (k\i.north);
      \draw (k\i.south) -- (0,-1.45);
      \draw (-0.575,-1.45) -- (0.575,-1.45);
      \draw (m\i.north) -- (-0.575,-1.45);
      \draw (n\i.north) -- (0.575,-1.45);
      \draw (n\i.south) -- (0.575,-3.3);
    \end{scope}
  }
  \draw[line width=1.4pt] (-1.25,-3.3) -- (9.5,-3.3);
  \node[font=\footnotesize, anchor=north west] at (-1.25,-3.35) {\iflangja{ネットワーク}{network}};
  % local access of core 0 to its own memory
  \draw[->, thick, ln-tip] (-0.12,-1.05) -- (-0.12,-1.33) -- (-0.42,-1.33) -- (-0.42,-1.95);
  \node[font=\scriptsize, text=ln-tip, anchor=east, inner sep=1pt] at (-0.64,-1.7)
    {\iflangja{高速}{fast}};
  % remote access from core 0 to the memory of node 2, through both NICs
  \begin{scope}[dashed, thick, ln-caution]
    \draw (0.12,-1.05) -- (0.12,-1.33) -- (0.8,-1.33) -- (0.8,-1.95);
    \draw (0.8,-2.45) -- (0.8,-3.05) -- (6.3,-3.05) -- (6.3,-2.45);
    \draw[->] (6.3,-1.95) -- (6.3,-1.33) -- (5.2,-1.33) -- (5.2,-1.95);
  \end{scope}
  \node[font=\scriptsize, text=ln-caution, anchor=north] at (3.4,-3.35)
    {\iflangja{他ノードのメモリ：ネットワーク経由（低速）}{memory of another node: via the network (slow)}};
\end{tikzpicture}

  \caption{分散共有メモリ: 各ノードは，コア，そのキャッシュ，ローカル
    メモリ，ネットワークインタフェースカード（NIC）からなる．}
  \label{fig:l18-distributed}
\end{figure}

コアとそのメモリおよびネットワークインタフェースは，それ自体ほとんど
1台の計算機であり，\emph{ノード}と呼ばれる．そのためこの種のマシンは
\emph{クラスタ}あるいは\emph{マルチコンピュータ}とも呼ばれる．コアは
自身のローカルメモリには単一プロセッサと同じ速さで到達できるが，他の
ノードが保持するデータはネットワークを渡らなければならず，はるかに長い
時間がかかる．アクセス時間はデータの所在に依存するので，このマシンは
\term{NUMA}[non-uniform memory access] である．

分散共有メモリは非常に多数のコアまでスケールし，クラスタやスーパー
コンピュータはこの方式で作られる．その理由は，ネットワークが速いから
ではない．ネットワークはメモリバスよりはるかに遅い．そうではなく，すべて
のコアが共有するメモリもバスも存在しないこと，そして，リモートアクセスを
すべて明示的に要求しなければならないプログラマが，通信について考えること
を強いられ，結果としてそれを最小限に保つことが理由である．\margin{目安と
して，ローカルの DRAM へは \SI{80}{\nano\second}，他方のソケットの
メモリへは \SI{130}{\nano\second}，InfiniBand で他のノードへは
1--\SI{2}{\micro\second} である．}

リモートアクセスはローカルアクセスよりはるかに高価なので，NUMA マシンの
性能はデータの配置で決まる．オペレーティングシステムは，各ページを，
それを最も多くアクセスするコアのノードのメモリに---スレッドのスタックの
ようなスレッド固有のデータであれば，そのスレッドが走るノードに---置き，
スレッドをそのノードに留めるべきである．そうすればメモリアクセスの
大半がローカルに収まる．ただし，複数のコアが使うページは，そのすべてに
とってローカルにはなりえない．\tip{NUMA マシンでは，各スレッドを，それ
が最も多く使うデータを保持するノードに留めること．Linux はページを，
最初にそれに触れたスレッドのノードに置く．\texttt{numactl} でこの選択を
上書きできる．}

\begin{note}[用語について]
  講義に従い，ここでは「分散共有メモリ」を，メモリがノードに分割されて
  いるすべてのマルチプロセッサに対して用いる．ネットワーク経由の明示的な
  要求で通信する場合，実際にはどのメモリも共有されていないが，本章では
  その場合もこの名前で呼ぶ．
  H\&P はこの名前を，ノードが1つの物理アドレス空間を実際に共有する
  マシン---どのコアも任意のアドレスを発行でき，相互結合網がリモート
  アドレスへのアクセスを，それを保持するノードへ運ぶ---に限って用い，
  明示的なメッセージだけで通信するマシンはクラスタと呼ぶ．講義はどちらも
  NUMA と呼ぶが，H\&P の用語では NUMA は，アドレス空間を共有し，アクセス
  時間がデータの所在に依存するマシンだけを指す．
\end{note}

\section{メッセージパッシングと共有メモリ}
\label{sec:l18-programs}

並列プログラムの書き方には，根本的に異なる2つの方法がある．
\source{Lesson 18，動画 11--14；H\&P §5.1．}

\begin{definition}[メッセージパッシング]
  \term[めっせーじぱっしんぐ]{メッセージパッシング}[message passing]の
  プログラミングモデルでは，各プロセスが自身のメモリを持ち，他のどの
  プロセスもそれにアクセスできない．プロセス間のデータ交換は，明示的に
  メッセージを送受信することによってのみ行われる．
\end{definition}

\begin{definition}[共有メモリ]
  \term[きょうゆうめもり]{共有メモリ}[shared memory]のプログラミング
  モデルでは，すべてのスレッドが1つのアドレス空間を共有し，他のスレッド
  が読む変数に書き込むだけでデータを交換する．
\end{definition}

以下では，同じ課題---$N$ 要素の配列の総和を $P$ 個のプロセスまたは
スレッドで求める---について2つのモデルを比較する．どちらのプログラムでも
各プロセスまたはスレッドは同じコードを実行し，0 から $P-1$ までの値を
とる識別子 \texttt{myPID} だけが異なる．

\subsection{メッセージパッシングのプログラム}

各プロセスは，自身に割り当てられた $N/P$ 要素だけを保持する．配列の各部分
を配布すること自体が明示的な作業であり（ここには示さない），最初にデータ
を持つプロセスが，他の各プロセスにその担当部分を送らなければならない．
次に各プロセスは自身の部分の総和を求め，プロセス 0 が他のすべての
プロセスの部分和を受け取って合算する（\cref{ex:l18-mp}）．

\needspace{8\baselineskip}
\begin{example}[メッセージパッシング]
  \label{ex:l18-mp}
  各プロセスは，自身の $N/P$ 要素をすでに \texttt{part} に持った状態で，
  次のコードを実行する．
  \begin{lstlisting}[language=C, numbers=none]
#define N 4096            /* 配列サイズ */
#define P 8               /* プロセス数 */
double part[N/P];         /* 自分の担当部分 */
double mySum = 0;
for (int i = 0; i < N/P; i++)
  mySum += part[i];
if (myPID == 0) {
  for (int p = 1; p < P; p++) {
    double pSum;
    recv(p, &pSum);       /* p の部分和を待つ */
    mySum += pSum;
  }
  printf("sum: %f\n", mySum);
} else
  send(0, &mySum);        /* プロセス 0 へ */
\end{lstlisting}
\end{example}

\subsection{共有メモリのプログラム}

共有メモリでは配列全体がすべてのスレッドから見えるので，\margin{実際には
ライブラリが用いられる．メッセージパッシングには MPI
（\texttt{MPI\_Reduce} がまさにこの総和である），共有メモリには OpenMP や
pthreads である．}%
配布すべきものは
何もない．各スレッドは，自身が担当するインデックスの範囲を計算し，その
部分和を共有の総和に加える（\cref{ex:l18-sm}）．\texttt{shared} と宣言
された変数は全スレッドに対して1つだけ存在し，それ以外の変数は各スレッド
固有である．

\begin{example}[共有メモリ]
  \label{ex:l18-sm}
  各スレッドは，共有された配列に対して次のコードを実行する．\caution{%
  \texttt{shared} は記法であって C の構文ではない．pthreads では大域変数
  とヒープが共有され，スタックは各スレッド固有である．}
  \begin{lstlisting}[language=C, numbers=none, morekeywords={shared}]
#define N 4096
#define P 8
shared double array[N];
shared double allSum = 0;
shared lock_t sumLock;
double mySum = 0;         /* スレッド固有 */
int lo = myPID * N / P;
int hi = (myPID + 1) * N / P;
for (int i = lo; i < hi; i++)
  mySum += array[i];
lock(&sumLock);
allSum += mySum;
unlock(&sumLock);
if (myPID == 0)           /* 本文を参照 */
  printf("sum: %f\n", allSum);
\end{lstlisting}
\end{example}

ロックは不可欠である．文 \texttt{allSum += mySum} は \texttt{allSum} を
読み，加算し，結果を書き戻す．2つのスレッドが，どちらかが書き込む前に
ともに古い値を読んでしまうと，2つの部分和の一方が失われる．ロックは，
この更新に一度に1つのスレッドだけを入れる．\tip{$P$ 個の部分和を順に
合算すると $P-1$ 段かかるが，木にすれば $\log_{2}P$ 段で済む．}%
またこのプログラムはまだ
未完成でもある．スレッド 0 は，他のすべてのスレッドが自身の和を加え
終わるまで待ってから合計を印字しなければならない．ロックと，この種の
待ち合わせは第20章の主題である．

\subsection{比較}

\cref{tab:l18-mp-sm} は2つのモデルの違いをまとめたものである．
メッセージパッシングでは，データの転送も，どのプロセスがどのデータを
持つかの決定も，すべてプログラマが書く．ハードウェアは，それ以外は独立な
ノードをネットワークで接続しさえすればよい．共有メモリでは，通信もデータ
の分配も共有変数を通じて暗黙に行われるが，ハードウェアは共有アドレス空間
を提供しなければならず，そこには全コアのキャッシュ内容の一貫性（第19章）
と，同期のための支援（第20章）が含まれる．2つのモデルは，明示的なコードを
必要とする場所も異なる．メッセージパッシングにはデータを配布するコードが
必要だが，受信そのものがデータの到着まで待ち合わせる．共有メモリでは配布の
コードは不要だが，あるスレッドが他のスレッドを待たなければならない箇所
すべてに明示的な同期が必要になる．\caution{受信が待つことは罠でもある．
2つのプロセスが互いに，受信より先に送信すると，どちらも永久に止まり
うる．}

\begin{table}[htbp]
  \centering\small
  \caption{メッセージパッシングと共有メモリの比較．}
  \label{tab:l18-mp-sm}
  \begin{tabular}{@{}lll@{}}
    \toprule
                             & メッセージパッシング & 共有メモリ \\
    \midrule
    通信                     & プログラマによる     & 自動 \\
    データの分配             & 手動                 & 自動 \\
    ハードウェアの支援       & 単純                 & 大がかり \\
    正しさのためのプログラミング & 難しい            & あまり難しくない \\
    性能のためのプログラミング   & 難しい            & 非常に難しい \\
    \bottomrule
  \end{tabular}
\end{table}

プログラマにとって，2つのモデルは難しさの配分が異なる．共有メモリの
プログラムは正しく書くのが容易である．逐次プログラムをループ単位で並列化
し，共有データの更新をロックで囲めばよい．一方で，速くするのははるかに
難しい．共有のコストがソースコードに現れないからである．ロックを待つ
スレッド，異なるコアのキャッシュ間を移動するデータ，NUMA マシンでの
リモートアクセスといったものである．メッセージパッシングのプログラムは
正しく書くのが難しい．データの分配と，すべての送信およびそれに対応する受信を手で
書かなければならないからである．しかし，いったん正しくなれば通信はすべて
コードに見えており，速くするための作業の多くはその過程で済んでいる．

\begin{keypoint}
  メッセージパッシングは通信を明示的にする．正しくプログラムするのは
  難しいが，そのコストは目に見える．共有メモリは通信を暗黙にする．正しく
  プログラムするのは容易だが，速くするのははるかに難しい．
\end{keypoint}

\section{共有メモリのハードウェア}
\label{sec:l18-mt-hardware}

共有メモリのプログラムのスレッドを，ハードウェアが実行する方法はいくつか
ある．\source{Lesson 18，動画 15；H\&P 第 3 章，§5.1．}
\begin{enumerate}
  \item \textbf{物理アドレス空間を共有する複数のコア．} どのコアも任意の
    物理アドレスを発行できる．メモリは集中していても（UMA），ノードに
    分散したうえで1つの物理アドレス空間を共有していてもよい（NUMA．
    \cref{sec:l18-distributed} の注を参照）．各スレッドはそれぞれのコアで
    走り，スレッドは本当に同時に実行される．
  \item \textbf{1つのコアによる時分割．} オペレーティングシステムが1つの
    スレッドをしばらく走らせ，その状態（PC とレジスタ）を保存して別の
    スレッドの状態を復元する．この\emph{コンテキストスイッチ}に特別な
    ハードウェアは要らず，スレッドは自動的に同じ物理メモリを使うが，ある
    時点で実行されるスレッドは1つだけであり，時分割では性能は何も得られ
    ない．
  \item \textbf{1つのコアによるハードウェアマルチスレッディング．} コア
    自身が複数のスレッドの状態を保持し，オペレーティングシステムを介さず
    にそれらを切り替える．コアがスレッドを切り替える頻度によって3つの
    変種が区別され，後のものほど多くのハードウェアの支援を必要とする．
    \begin{itemize}
      \item \term[そりゅうどまるちすれっでぃんぐ]{粗粒度マルチスレッディング}[coarse-grained multithreading]%
        は，数サイクル以上経ってから別のスレッドに切り替える．典型的には，
        走っているスレッドが，最終レベルキャッシュのミスのような遅い事象
        を待たなければならなくなったときである．
      \item \term[さいりゅうどまるちすれっでぃんぐ]{細粒度マルチスレッディング}[fine-grained multithreading]%
        は，毎サイクル別のスレッドに切り替えることができる．
      \item \term[どうじまるちすれっでぃんぐ]{同時マルチスレッディング}[simultaneous multithreading]%
        \index{SMT|see{同時マルチスレッディング}}（SMT）は，同じサイクル
        に複数のスレッドの命令を実行する．Intel の Hyper-Threading は SMT
        の実装である．
    \end{itemize}
\end{enumerate}

\section{マルチスレッディングの性能}
\label{sec:l18-mt-performance}

2つのスレッドと，サイクル当たり4命令を扱えるスーパースカラコアを考える．
\source{Lesson 18，動画 16--17；H\&P 第 3 章；
\cite{olukotun1996,tullsen1995}．}1つのスレッドが1サイクルの4スロットを
すべて埋めることは稀である．命令間の依存，予測を外した分岐，キャッシュ
ミスによって多くのスロットが空いたままになり，まったく進めないサイクルも
ある．\cref{fig:l18-mt-timing} は，4つの構成がスロットをどう使うかを
模式的に示す．4つのいずれにおいても，空いたスロットを得られなかった
命令は次のサイクルまで待ち，キャッシュミスのような事象を待つスレッドは，コアを保持
しているかどうかによらず同じサイクル数だけ待つ．

\begin{figure}[htbp]
  \centering
  % Schematic use of the four slots per cycle of a core by two threads over
% ten cycles: (a) time-sharing, (b) chip multiprocessor, (c) fine-grained
% multithreading, (d) simultaneous multithreading.
% Alone, thread 1 executes 3,2,3,-,-,1,1,-,4,2 and thread 2
% 3,1,2,-,2,-,-,4,1,3 instructions in the ten cycles (16 each; "-" is a
% cycle spent waiting for an event such as a cache miss).  One model for all
% panels: instructions for which no slot is free wait for the next cycle
% (with SMT, 4 slots are split 2+2 when both threads have more ready), and
% a wait elapses whether or not the thread holds the core.  (a) switches
% threads after five cycles; (c) alternates between the threads that are
% not waiting.
\begin{tikzpicture}[x=1cm,y=1cm, font=\scriptsize\sffamily]
  \def\lncw{0.24}
  % \lnslots{x0}{y0}{cycle/slots of thread 1/slots of thread 2, ...}
  \def\lnslots#1#2#3{%
    \foreach \c/\a/\b in {#3} {
      \foreach \s in {1,...,4} {
        \pgfmathtruncatemacro\lnk{\s<=\a ? 1 : (\s<=\a+\b ? 2 : 0)}
        \ifcase\lnk \def\lnf{white}\or \def\lnf{ln-accent!60}\or \def\lnf{ln-tip!45}\fi
        \draw[ln-rule, fill=\lnf] ({#1+(\c-1)*\lncw},{#2-(\s-1)*\lncw})
          rectangle ++(\lncw,-\lncw);
      }
    }
    \draw ({#1},{#2}) rectangle ++({10*\lncw},{-4*\lncw});
  }
  % (a) time-sharing: thread 1 for five cycles, then thread 2
  \node[anchor=south west, inner sep=1pt] at (0,0.05)
    {\iflangja{(a) 時分割}{(a) time-sharing}};
  \lnslots{0}{0}{1/3/0,2/2/0,3/3/0,4/0/0,5/0/0,6/0/3,7/0/1,8/0/2,9/0/0,10/0/2}
  \node[anchor=north, text=ln-muted] at (1.2,-0.98) {\iflangja{コア：16 命令}{core: 16 instr.}};
  % (b) CMP: one core per thread
  \node[anchor=south west, inner sep=1pt] at (3.0,0.05) {(b) CMP};
  \lnslots{3.0}{0}{1/3/0,2/2/0,3/3/0,4/0/0,5/0/0,6/1/0,7/1/0,8/0/0,9/4/0,10/2/0}
  \node[anchor=north, text=ln-muted] at (4.2,-0.98) {\iflangja{コア 1：16 命令}{core 1: 16 instr.}};
  \lnslots{3.0}{-1.44}{1/0/3,2/0/1,3/0/2,4/0/0,5/0/2,6/0/0,7/0/0,8/0/4,9/0/1,10/0/3}
  \node[anchor=north, text=ln-muted] at (4.2,-2.42) {\iflangja{コア 2：16 命令}{core 2: 16 instr.}};
  % (c) fine-grained multithreading: one thread per cycle
  \node[anchor=south west, inner sep=1pt] at (6.0,0.05)
    {\iflangja{(c) 細粒度 MT}{(c) fine-grained MT}};
  \lnslots{6.0}{0}{1/3/0,2/0/3,3/2/0,4/0/1,5/3/0,6/0/2,7/0/0,8/1/0,9/0/2,10/1/0}
  \node[anchor=north, text=ln-muted] at (7.2,-0.98) {\iflangja{コア：18 命令}{core: 18 instr.}};
  % (d) SMT: both threads in the same cycle
  \node[anchor=south west, inner sep=1pt] at (9.0,0.05) {(d) SMT};
  \lnslots{9.0}{0}{1/2/2,2/1/1,3/2/1,4/2/2,5/1/0,6/0/2,7/0/0,8/1/0,9/1/3,10/0/1}
  \node[anchor=north, text=ln-muted] at (10.2,-0.98) {\iflangja{コア：22 命令}{core: 22 instr.}};
  % axis hints
  \draw[->, ln-muted] (6.0,-1.7) -- (8.4,-1.7)
    node[midway, below, inner sep=1.5pt] {\iflangja{サイクル}{cycles}};
  \draw[->, ln-muted] (-0.12,0) -- (-0.12,-0.96)
    node[midway, rotate=90, above, inner sep=1.5pt] {\iflangja{スロット}{slots}};
  % legend
  \begin{scope}[xshift=9.0cm, yshift=-1.6cm]
    \draw[ln-rule, fill=ln-accent!60] (0,0) rectangle ++(0.24,-0.24);
    \node[anchor=west] at (0.3,-0.125) {\iflangja{スレッド 1}{thread 1}};
    \draw[ln-rule, fill=ln-tip!45] (0,-0.4) rectangle ++(0.24,-0.24);
    \node[anchor=west] at (0.3,-0.525) {\iflangja{スレッド 2}{thread 2}};
    \draw[ln-rule, fill=white] (0,-0.8) rectangle ++(0.24,-0.24);
    \node[anchor=west] at (0.3,-0.925) {\iflangja{空きスロット}{empty slot}};
  \end{scope}
\end{tikzpicture}

  \caption{2つのスレッドによる，コアのサイクル当たり4スロットの
    10サイクルにわたる使われ方の模式図．各列が1サイクル，各行が1
    スロットである．単独ではどちらのスレッドもこの10サイクルで16命令を
    実行する (b)．(a) は5サイクルごとにスレッドを切り替え，(c) は
    待っていないスレッドの間で交互に切り替える．}
  \label{fig:l18-mt-timing}
\end{figure}

\begin{description}
  \item[時分割] 1つのコアが一度に1つのスレッドを走らせる．2つのスレッド
    を合わせても，1つのスレッドより多くの仕事は完了しない．コンテキスト
    スイッチのたびに時間がかかるので（図には描いていない），実際には
    わずかに少ない．コストはコア1個分である．
  \item[チップマルチプロセッサ] 各スレッドがそれぞれのコアで走る．この
    ように使われるマルチコアプロセッサを
    \term[ちっぷまるちぷろせっさ]{チップマルチプロセッサ}[chip multiprocessor]%
    \index{CMP|see{チップマルチプロセッサ}}（CMP）という．どちらの
    スレッドも全速で走るのでスループットは2倍になるが，チップ面積と電力も
    2倍になり，しかも各コアは単一スレッドの場合と同じだけスロットを空けた
    ままにする．
  \item[細粒度マルチスレッディング] スレッドごとに別のレジスタ組を持つ
    1つのコアが，サイクルごとにスレッドを切り替えるので，一方のスレッドが
    待たなければならないサイクルを他方が使う．コアはより忙しくなり，追加
    のハードウェアはスレッドごとの小さな状態だけで，コアのコストのおよそ
    \SI{5}{\percent} である．ただしサイクル内の空きスロットは空いたままで
    ある．
  \item[SMT] 1つのコアが同じサイクルに両方のスレッドの命令を混ぜ，単一の
    スレッドなら空いたままになるスロットの多くを埋める．\sidenote{SMT は
    1995 年に提案された\cite{tullsen1995}．Intel は 2002 年の Pentium~4
    にハイパースレッディングとして搭載した．コア当たり2スレッドというこの
    数は x86 では今も普通である．一方 IBM の POWER8 と POWER9 は，
    1つのコアで最大8スレッドを走らせる．Intel は 2024 年のクライアント
    向けチップの Lion Cove コアで SMT を再び取りやめ，その面積はコア数に
    充てる方が良いと判断した．}%
    スループットは
    細粒度マルチスレッディングより高くなるが，スレッドが1組の資源を共有
    するため，2個のコアには及ばない．SMT は細粒度マルチスレッディングより
    多くの追加ハードウェアを必要とし，コアのコストのおよそ
    \SI{10}{\percent} である（\cref{sec:l18-smt-hw}）．
\end{description}

SMT がどれだけの効果を上げるかはスレッドに依存する．両者の命令が同じ実行
ユニットを奪い合う場合，一方のスレッドの命令は他方が空けたスロットを
ほとんど見つけられず，得られる利得は小さい．2つのスレッドの資源要求が
互いに補い合うときにのみ，SMT は CMP のスループットに近づく．

\begin{example}
  \label{ex:l18-smt-cmp}
  SMT がコアのコストを \SI{10}{\percent} 増やし，2つのスレッドに対する
  スループットを \SI{35}{\percent} 高めるとする．SMT なしの2コアの CMP は，
  スループットもコストも2倍にする．スレッドを時分割する1個のコアを基準に
  すると，SMT のコアはコスト当たり $1.35/1.10\approx 1.23$ 倍の，CMP は
  $2/2=1$ 倍のスループットを与える．ただし絶対的なスループットは CMP の
  方が高い．2つのスレッドが同じ実行ユニットを奪い合い，SMT が
  スループットをわずか \SI{5}{\percent} しか高めないのであれば，SMT の
  コアは $1.05/1.10\approx 0.95$ となり，2個のコアの方がハードウェアの
  使い方として優れることになる．
\end{example}

\begin{keypoint}
  CMP は最も高いスループットを2倍のコストで与える．SMT はおよそ
  \SI{10}{\percent} の追加コストで，より小さな利得を与える．スレッドの
  資源要求が互いに補い合うときは SMT の方がハードウェアの効率的な使い方
  であり，スレッドが同じ資源を奪い合うときは2個のコアの方が良い選択に
  なりうる．
\end{keypoint}

\section{SMT のためのハードウェアの変更}
\label{sec:l18-smt-hw}

SMT は，第7章・第8章のアウトオブオーダコアにわずかな変更を加えるだけで
実現できる（\cref{fig:l18-smt-hw}）．\source{Lesson 18，動画 18；H\&P
第 3 章；\cite{tullsen1995}．}

\begin{figure}[htbp]
  \centering
  % An out-of-order core with SMT support and separate ROBs: per-thread PC,
% RAT, ROB and architectural registers; everything else is shared.
\begin{tikzpicture}[x=1cm,y=1cm, font=\footnotesize\sffamily, >={Stealth[length=1.8mm]},
  sh/.style={draw, fill=ln-box, minimum height=0.75cm, minimum width=1.5cm, inner sep=2pt, align=center},
  pt/.style={draw, fill=ln-accent!15, minimum height=0.36cm, minimum width=1.2cm, inner sep=1pt, font=\scriptsize\sffamily}]
  % ---- front end
  \node[pt, minimum width=0.9cm] (pc1) at (0,0.22) {PC 1};
  \node[pt, minimum width=0.9cm] (pc2) at (0,-0.22) {PC 2};
  \draw[fill=white] (0.85,0.45) -- (1.15,0.3) -- (1.15,-0.3) -- (0.85,-0.45) -- cycle;
  \draw[->] (pc1.east) -- (0.85,0.22);
  \draw[->] (pc2.east) -- (0.85,-0.22);
  \node[sh] (fetch) at (2.3,0) {\iflangja{フェッチ}{fetch}};
  \node[sh] (dec) at (4.4,0) {\iflangja{デコード}{decode}};
  \draw[->] (1.15,0) -- (fetch.west);
  \draw[->] (fetch.east) -- (dec.west);
  \node[pt] (rat1) at (6.6,0.22) {RAT 1};
  \node[pt] (rat2) at (6.6,-0.22) {RAT 2};
  \node[draw, fit=(rat1)(rat2), inner sep=2pt] (ren) {};
  \node[anchor=south, font=\scriptsize] at (ren.north) {\iflangja{リネーム}{rename}};
  \draw[->] (dec.east) -- (ren.west);
  % ---- back end
  \node[sh] (rs) at (1.1,-1.9) {RS};
  \node[sh] (ex) at (3.2,-1.9) {\iflangja{実行}{execute}};
  \node[pt, minimum width=1.6cm] (rob1) at (5.6,-1.68) {ROB 1};
  \node[pt, minimum width=1.6cm] (rob2) at (5.6,-2.12) {ROB 2};
  \node[draw, fit=(rob1)(rob2), inner sep=2pt] (robs) {};
  \node[pt] (arf1) at (8.0,-1.68) {\iflangja{レジスタ 1}{ARF 1}};
  \node[pt] (arf2) at (8.0,-2.12) {\iflangja{レジスタ 2}{ARF 2}};
  \draw[->] (ren.south) -- ++(0,-0.35) -| (rs.north);
  \draw[->] (ren.south) -- ++(0,-0.35) -| (robs.north);
  \draw[->] (rs.east) -- (ex.west);
  \draw[->] (ex.east) -- (robs.west |- ex.east) node[midway, above, font=\scriptsize] {CDB};
  \draw[->] (ex.south) -- ++(0,-0.4) -| (rs.south);
  \draw[->] (robs.east |- rob1.east) -- (arf1.west);
  \draw[->] (robs.east |- rob2.east) -- (arf2.west);
  \node[anchor=north, font=\scriptsize] at (6.95,-2.3) {\iflangja{コミット}{commit}};
  % ---- legend
  \draw[fill=ln-accent!15] (0,-3.05) rectangle ++(0.34,-0.25);
  \node[anchor=west, font=\scriptsize] at (0.4,-3.175) {\iflangja{スレッドごと}{one per thread}};
  \draw[fill=ln-box] (2.6,-3.05) rectangle ++(0.34,-0.25);
  \node[anchor=west, font=\scriptsize] at (3.0,-3.175) {\iflangja{スレッド間で共有}{shared by the threads}};
\end{tikzpicture}

  \caption{2つのスレッドのための SMT を備え，ROB を分離したアウトオブ
    オーダコア．ROB を共有する場合，スレッドごとに保持されるのは
    プログラムカウンタ，RAT，アーキテクチャレジスタだけである．}
  \label{fig:l18-smt-hw}
\end{figure}

\begin{description}
  \item[フェッチ] 各スレッドが自身の PC を必要とする．フェッチステージは
    両方の PC から命令を取り，フェッチした各命令に，それが属するスレッド
    の印を付ける．
  \item[デコード] デコードはスレッドに依存せず，変更はない．
  \item[リネーム] リネームの論理は変わらないが，各スレッドが自身の RAT を
    必要とする．スレッド 1 のレジスタ R1 とスレッド 2 のレジスタ R1 は，
    別のアーキテクチャレジスタだからである．
  \item[ディスパッチと実行] 変更はない．リザベーションステーション，実行
    ユニット，CDB は共有される．リネームの後，オペランドは両方のスレッド
    にわたって一意なタグで識別されるので，異なるスレッドの命令が混同され
    ることはない．
  \item[ROB] 各スレッドが，自身のコミット点を持つ別々の ROB を持つことも
    できる．しかし ROB は大きな構造なので，複製するのは高価である．実際
    にはより多くの場合，両方のスレッドの命令が1つの共有 ROB に混在して並
    び，各エントリがその命令のスレッドを記録する．\margin{Pentium~4 は共有
    ではなく分割した．両方のスレッドが走っている間，各スレッドは 126 個の
    ROB エントリの半分を得た\cite{marr2002}．}%
    その小さな代償は，
    あるスレッドの完了した命令が，他方のスレッドのより古い命令がコミット
    するのを待たされることがある点である．
  \item[コミット] 各スレッドが自身のアーキテクチャレジスタファイル（ARF）
    を必要とし，そのスレッドの命令はそこにコミットする．
\end{description}

ROB を共有する場合，複製される構造---PC，RAT，ARF---はいずれも小さく，
コアの高価な部分---ROB，リザベーションステーション，実行ユニット，
キャッシュ，分岐予測器---は共有される．SMT がコアのコストをおよそ
\SI{10}{\percent} しか増やさないのは，このためである．\margin{Intel は
ハイパースレッディングの Pentium~4 での面積コストをダイのおよそ
\SI{5}{\percent}，スループットの向上を \SI{15}{\percent}--%
\SI{30}{\percent} としている\cite{marr2002}．}

\section{SMT とデータキャッシュ・TLB}
\label{sec:l18-smt-tlb}

SMT のコアを共有するスレッドは，それぞれ自身の仮想アドレス空間を持つ別の
プロセスに属することがある．\source{Lesson 18，動画 19；H\&P 付録 B．}%
そのとき同じ仮想アドレスはスレッドごとに異なるデータを指すので，
第13章・第14章のキャッシュと TLB は，両者を混同してはならない．
\begin{itemize}
  \item \term*{仮想アドレスキャッシュ}は，仮想アドレスだけでインデックス
    し，タグも付ける．一方のスレッドのためにキャッシュされたブロックは，
    他方のスレッドの同じ仮想アドレスへのアクセスに対してヒットし，誤った
    データを返す．この種のキャッシュを SMT で使えるのは，各エントリが
    スレッドも記録し，ヒットにスレッドの一致も要求する場合に限られる．
  \item \term*{VIPT キャッシュ}と物理アドレスキャッシュは，物理タグを
    比較する．2つのスレッドが同じブロックにヒットするのは，同じ物理
    アドレスにアクセスしたときだけであり，そのときは実際に同じデータに
    アクセスしている．これらのキャッシュは，TLB が各スレッドに正しい変換
    を与える限り，SMT でも正しく動作する．
  \item したがって \term*{TLB} は，スレッドを意識しなければならない．
    各エントリはそれが属するスレッドを記録し，アクセスのスレッド識別子が
    探索に渡され，ヒットにはスレッドと仮想ページ番号の両方の一致が要求
    される（\cref{fig:l18-smt-tlb}）．\caution{区別すべきはスレッドでは
    なくアドレス空間である．同じプロセスの2つのスレッドは変換を共有する．
    実際の TLB は各エントリにアドレス空間識別子---x86-64 では PCID，
    Arm では ASID---を付ける．}
\end{itemize}

\begin{figure}[htbp]
  \centering
  % Thread-aware TLB of an SMT core: every entry records the thread it
% belongs to, and a lookup must match both the thread and the VPN.
\begin{tikzpicture}[x=1cm,y=1cm, font=\footnotesize\sffamily, >={Stealth[length=1.8mm]},
  cell/.style={draw, minimum height=0.42cm, inner sep=1pt, font=\footnotesize\ttfamily},
  hd/.style={inner sep=1pt, font=\scriptsize\sffamily, text=ln-muted}]
  % ---- lookup key
  \node[cell, fill=ln-accent!15, minimum width=0.8cm] (kt) at (0,0.21) {1};
  \node[cell, fill=ln-accent!15, minimum width=1.2cm, anchor=west] (kv) at (kt.east) {0x12};
  \node[hd, anchor=south] at (kt.north) {\iflangja{スレッド}{thread}};
  \node[hd, anchor=south] at (kv.north) {VPN};
  % ---- TLB entries: thread | VPN | frame
  \foreach \r/\t/\v/\f in {0/0/0x12/0x3A, 1/1/0x12/0x07, 2/0/0x40/0x15, 3/1/0x2C/0x51} {
    \ifnum\r=1 \def\lnfill{ln-tip!25}\else \def\lnfill{white}\fi
    \node[cell, fill=\lnfill, minimum width=0.8cm] (t\r) at (4.2,{0.63-\r*0.42}) {\t};
    \node[cell, fill=\lnfill, minimum width=1.2cm, anchor=west] (v\r) at (t\r.east) {\v};
    \node[cell, fill=\lnfill, minimum width=1.2cm, anchor=west] (f\r) at (v\r.east) {\f};
  }
  \node[hd, anchor=south] at (t0.north) {\iflangja{スレッド}{thread}};
  \node[hd, anchor=south] at (v0.north) {VPN};
  \node[hd, anchor=south] at (f0.north) {\iflangja{フレーム}{frame}};
  \node[font=\footnotesize, anchor=north] at (v3.south) {TLB};
  % ---- lookup and result
  \draw[->] (kv.east) -- (t1.west |- kv.east)
    node[midway, above, font=\scriptsize, align=center]
    {\iflangja{スレッドと VPN が}{thread and VPN}\\\iflangja{ともに一致}{must both match}};
  \node[cell, fill=ln-tip!25, minimum width=1.2cm] (res) at (9.2,0.21) {0x07};
  \node[hd, anchor=south] at (res.north) {\iflangja{物理フレーム}{physical frame}};
  \draw[->] (f1.east) -- (res.west);
  \node[hd, anchor=north west, align=left] at (3.8,-1.35)
    {\iflangja{同じ VPN 0x12 でもスレッド 0 はフレーム 0x3A に対応する}{VPN 0x12 of thread 0 maps to frame 0x3A instead}};
\end{tikzpicture}

  \caption{スレッドを意識した TLB．同じ仮想ページ番号が，2つのスレッドに
    対して別々のフレームに変換される．}
  \label{fig:l18-smt-tlb}
\end{figure}

\section{SMT とキャッシュ性能}
\label{sec:l18-smt-cache}

コアのキャッシュは，そのコアのすべての SMT スレッドに共有される．これに
は1つの利点と，1つの深刻な欠点がある．\source{Lesson 18，動画 20--22；
H\&P 第 3 章．}
\begin{itemize}
  \item \textbf{高速なデータ共有．} 一方のスレッドが格納した値を他方の
    スレッドがロードするとき，そのロードはデータをすでにキャッシュ内に
    見つける．
  \item \textbf{容量の共有．} スレッドはキャッシュの容量と連想度も共有
    しており，それらを奪い合う．
\end{itemize}

この競合は，スレッドの\term[わーきんぐせっと]{ワーキングセット}[working set]%
---ある時点でそのスレッドが実際に使っているブロックの集合---によって
記述される．2つの SMT スレッドのワーキングセットの大きさを
$\mathit{WS}_0$，$\mathit{WS}_1$ とし，両者がともに使う部分の大きさを
$\mathit{WS}_{01}$ とする．
\begin{equation}
  \mathit{WS}_0 + \mathit{WS}_1 - \mathit{WS}_{01} \;\le\; \text{キャッシュ容量}
  \label{eq:l18-ws}
\end{equation}
が成り立つ限り，どちらのスレッドも，単独で走るときとほとんど変わらない
頻度で自身のブロックをキャッシュ内に見つける．合わせたワーキングセットが
キャッシュを超えると，スレッドは互いのブロックを追い出し続け，キャッシュ
は\emph{スラッシング}を起こす．それぞれのワーキングセットが単独なら
収まる場合，SMT はスレッドを1つずつ走らせるより遅くなりうる．時分割される
スレッドは数百万サイクルの間キャッシュを独占し，コンテキストスイッチの
たびに詰め直すだけであるのに対し，SMT のスレッドは絶えず互いのブロックを
追い出すからである．\caution{SMT がスループットを高めるのは，キャッシュ
がすべてのスレッドのワーキングセットを同時に保持できる場合だけである．}%
したがってスケジューラは，合わせたワーキングセットがキャッシュに収まる
スレッドどうしを一緒に走らせ，ワーキングセットだけでキャッシュを超える
スレッドは SMT の相手なしで走らせるべきである．そのようなスレッドは，
自身は何も得ないまま相手のブロックを追い出してしまうからである．

\begin{example}
  あるコアは \SI{64}{KB} の L1 データキャッシュを持つ．\margin{実際の L1
  データキャッシュはこれより小さい．AMD の Zen~4 で 32~KB，Intel の
  Golden Cove で 48~KB であり，その後ろにコア当たり 1--2~MB の L2 が
  ある．}%
  スレッド A，B，C
  のワーキングセットはそれぞれ \SI{40}{KB}，\SI{36}{KB}，\SI{28}{KB} で
  あり，いずれも単独ならキャッシュに収まる．A と B は共通のデータを扱って
  おり，そのうち \SI{20}{KB} が両方のワーキングセットに現れるので，SMT で
  一緒に走らせるのに必要なのは $40+36-20=\SI{56}{KB}$ であり，これは
  収まる．A と C はデータを共有せず，合わせて \SI{68}{KB} を必要とする．
  SMT ではスラッシングを起こし，時分割した方が速い可能性が高い．
\end{example}

\begin{keypoint}[章のまとめ]
  マルチプロセッサは MIMD マシンである．パイプラインを広げても収穫逓減と
  なり，周波数を上げると電力がおよそ $f^{3}$ で増えるのに対し，コアを
  複製すればムーアの法則によるトランジスタをスループットに変えられるため，
  マルチプロセッサはますます速くなる単一プロセッサに取って代わった---ただし
  それは並列プログラムに限られ，その高速化率は頭打ちになる
  （\cref{eq:l18-amdahl-cores}）．集中共有メモリ（UMA，SMP）は単純だが，
  メモリ容量と帯域によって少数のコアに制限される．分散共有メモリ（NUMA）
  は大規模なマシンまでスケールするが，リモートアクセスが遅いという代償を
  払う．メッセージパッシングは正しく書くのが難しく，共有メモリは速くする
  のが難しい．1つのコアは，時分割，粗粒度または細粒度マルチスレッディング，
  あるいは SMT によって複数のスレッドを走らせられる．SMT が複製するのは
  PC，RAT，ARF 程度であり，スレッドを意識した TLB を必要とし，共有された
  キャッシュが合わせたワーキングセットを保持できる場合にのみ引き合う
  （\cref{eq:l18-ws}）．メモリを共有するコアは，キャッシュの一貫性も
  保たなければならない．これが第19章の主題である．
\end{keypoint}

\end{document}
